% GENERATED FILE. Edit canonical lessons, then run npm run generate:books.
% canonical-source-hash: fnv1a64:92cd256f2825141f
% canonical-lessons: BN-C429-kalboishakhi, BN-C429-shilabrishti, BN-C429-gorjano, BN-C429-kapano, BN-C429-nuye-pora

\chapter{Storms and Hail}
\label{ch:bn-storms-and-hail}
\hlchaptermodality{429}

\begin{chapteropening}
\textbf{By the end of this chapter:} \emph{I can describe a spring storm, hail, thunder and trees bent by the wind.}

Complete the last lesson of chapter 429: I can describe a spring storm, hail, thunder and trees bent by the wind.
\end{chapteropening}

\section[\texorpdfstring{\bn{কালবৈশাখী} (kālboishākhī)}{কালবৈশাখী (kālboishākhī)}]{\bn{কালবৈশাখী} (kālboishākhī) — a nor'wester; a spring thunderstorm}

\label{lesson:BN-C429-kalboishakhi}



\begin{quote}
Before the new one: say the Bengali for damp (and dank), then the Bengali for to be covered over.
\end{quote}

\subsection*{You'll want to know: \bn{কালবৈশাখী}}
\textbf{\bn{কালবৈশাখী}} — \emph{kālboishākhī} — \textquotedblleft{}a nor'wester; a spring thunderstorm\textquotedblright{}.

\begin{quote}
\bn{বিকেলে} \bn{কালবৈশাখী} \bn{এসে} \bn{অনেক} \bn{গাছ} \bn{ভেঙে} \bn{দিল।}
\end{quote}

\begin{quote}
\emph{bikele kālboishākhī ese ônek gāchh bhenge dilô.}
\end{quote}

\begin{quote}
In the afternoon a nor'wester came and broke many trees.
\end{quote}

The sudden violent storms of April and May, in the month of Baishakh.

\begin{cousinweb}
\textbf{\bn{কালবৈশাখী}} (\emph{kālboishākhī}) joins \bn{কাল}, \textquotedblleft{}black, deadly\textquotedblright{}, to \bn{বৈশাখী}, \textquotedblleft{}of the month of Baishakh\textquotedblright{}, from Sanskrit \emph{vaiśākha}, the spring month (Monier-Williams): the black storm of Baishakh.
\end{cousinweb}

\begin{grammarlens}[title={What you've built}]
One more word for getting around.
\end{grammarlens}

\subsection*{Your turn}
\begin{itemize}
\raggedright
  \item \emph{Say it:} \emph{kālboishākhī}
  \item \emph{Say it:} \emph{kālboishākhī}, once more
  \item \emph{Say it:} \emph{bikele kālboishākhī ese ônek gāchh bhenge dilô.} — In the afternoon a nor'wester came and broke many trees.
  \item \emph{Recall:} say the Bengali for damp (and dank), then the Bengali for to be covered over, then say \emph{kālboishākhī} again
\end{itemize}

\begin{tcolorbox}[breakable,colback=teal!4,colframe=teal!35!black,title={Before you move on}]
What does \bn{কালবৈশাখী} mean? (\textquotedblleft{}A nor'wester; a spring thunderstorm\textquotedblright{}.) Say it once more.
\end{tcolorbox}
\section[\texorpdfstring{\bn{শিলাবৃষ্টি} (shilābrishṭi)}{শিলাবৃষ্টি (shilābrishṭi)}]{\bn{শিলাবৃষ্টি} (shilābrishṭi) — hail; a hailstorm}

\label{lesson:BN-C429-shilabrishti}



\begin{quote}
Before the new one: say the Bengali for to be covered over, then the Bengali for a nor'wester.
\end{quote}

\subsection*{You'll want to know: \bn{শিলাবৃষ্টি}}
\textbf{\bn{শিলাবৃষ্টি}} — \emph{shilābrishṭi} — \textquotedblleft{}hail; a hailstorm\textquotedblright{}.

\begin{quote}
\bn{ঝড়ের} \bn{সঙ্গে} \bn{শিলাবৃষ্টি} \bn{হল।}
\end{quote}

\begin{quote}
\emph{jhôṛer sôngge shilābrishṭi hôlô.}
\end{quote}

\begin{quote}
Hail came with the storm.
\end{quote}

\begin{cousinweb}
\textbf{\bn{শিলাবৃষ্টি}} (\emph{shilābrishṭi}) joins Sanskrit \emph{śilā}, \textquotedblleft{}a stone\textquotedblright{}, to \emph{vṛṣṭi}, \textquotedblleft{}rain\textquotedblright{} (Monier-Williams for both): a rain of stones.
\end{cousinweb}

\begin{grammarlens}[title={What you've built}]
One more word for getting around.
\end{grammarlens}

\subsection*{Your turn}
\begin{itemize}
\raggedright
  \item \emph{Say it:} \emph{shilābrishṭi}
  \item \emph{Say it:} \emph{shilābrishṭi}, once more
  \item \emph{Say it:} \emph{jhôṛer sôngge shilābrishṭi hôlô.} — Hail came with the storm.
  \item \emph{Recall:} say the Bengali for to be covered over, then the Bengali for a nor'wester, then say \emph{shilābrishṭi} again
\end{itemize}

\begin{tcolorbox}[breakable,colback=teal!4,colframe=teal!35!black,title={Before you move on}]
What does \bn{শিলাবৃষ্টি} mean? (\textquotedblleft{}Hail; a hailstorm\textquotedblright{}.) Say it once more.
\end{tcolorbox}
\section[\texorpdfstring{\bn{গর্জানো} (gôrjāno)}{গর্জানো (gôrjāno)}]{\bn{গর্জানো} (gôrjāno) — to roar; to rumble (of thunder or the sea)}

\label{lesson:BN-C429-gorjano}



\begin{quote}
Before the new one: say the Bengali for a nor'wester, then the Bengali for hail.
\end{quote}

\subsection*{You'll want to know: \bn{গর্জানো}}
\textbf{\bn{গর্জানো}} — \emph{gôrjāno} — \textquotedblleft{}to roar; to rumble (of thunder or the sea)\textquotedblright{}.

\begin{quote}
\bn{সারারাত} \bn{মেঘ} \bn{গর্জাল।}
\end{quote}

\begin{quote}
\emph{sārārāt megh gôrjālô.}
\end{quote}

\begin{quote}
The clouds rumbled all night.
\end{quote}

Past (he/she …): \bn{সেই} \bn{সমুদ্র} \bn{গর্জাল}, \emph{sei sômudrô gôrjālô}.

\begin{cousinweb}
\textbf{\bn{গর্জানো}} (\emph{gôrjāno}) probably goes back to Sanskrit \emph{gárjati}, \textquotedblleft{}thunders, roars\textquotedblright{} (CDIAL 4046), with the Bengali verb ending -\bn{আনো}; CDIAL lists no Bengali form there, so the link is hedged.
\end{cousinweb}

\begin{grammarlens}[title={What you've built}]
One more word for getting around.
\end{grammarlens}

\subsection*{Your turn}
\begin{itemize}
\raggedright
  \item \emph{Say it:} \emph{gôrjāno}
  \item \emph{Say it:} \emph{gôrjāno}, once more
  \item \emph{Say it:} \emph{sārārāt megh gôrjālô.} — The clouds rumbled all night.
  \item \emph{Recall:} say the Bengali for a nor'wester, then the Bengali for hail, then say \emph{gôrjāno} again
\end{itemize}

\begin{tcolorbox}[breakable,colback=teal!4,colframe=teal!35!black,title={Before you move on}]
What does \bn{গর্জানো} mean? (\textquotedblleft{}To roar; to rumble (of thunder or the sea)\textquotedblright{}.) Say it once more.
\end{tcolorbox}
\section[\texorpdfstring{\bn{কাঁপানো} (kā̃pāno)}{কাঁপানো (kā̃pāno)}]{\bn{কাঁপানো} (kā̃pāno) — to shake (something); to make tremble}

\label{lesson:BN-C429-kapano}



\begin{quote}
Before the new one: say the Bengali for hail, then the Bengali for to roar.
\end{quote}

\subsection*{You'll want to know: \bn{কাঁপানো}}
\textbf{\bn{কাঁপানো}} — \emph{kā̃pāno} — \textquotedblleft{}to shake (something); to make tremble\textquotedblright{}.

\begin{quote}
\bn{ঝড়} \bn{জানালাগুলো} \bn{কাঁপাল।}
\end{quote}

\begin{quote}
\emph{jhôṛ jānālāgulo kā̃pālô.}
\end{quote}

\begin{quote}
The storm rattled the windows.
\end{quote}

Past (he/she …): \bn{সেই} \bn{ঝড়} \bn{জানালা} \bn{কাঁপাল}, \emph{sei jhôṛ jānālā kā̃pālô}.

The causing partner of \bn{কাঁপা}, to shiver: you \bn{কাঁপো} with cold, and the wind \bn{কাঁপায়} the door.

\begin{cousinweb}
\textbf{\bn{কাঁপানো}} (\emph{kā̃pāno}) is the causing form of \textbf{\bn{কাঁপা}}, \textquotedblleft{}to shiver, to tremble\textquotedblright{}, from Sanskrit \emph{kampate}, \textquotedblleft{}trembles\textquotedblright{} (CDIAL 2767).
\end{cousinweb}

\begin{grammarlens}[title={What you've built}]
One more word for getting around.
\end{grammarlens}

\subsection*{Your turn}
\begin{itemize}
\raggedright
  \item \emph{Say it:} \emph{kā̃pāno}
  \item \emph{Say it:} \emph{kā̃pāno}, once more
  \item \emph{Say it:} \emph{jhôṛ jānālāgulo kā̃pālô.} — The storm rattled the windows.
  \item \emph{Recall:} say the Bengali for hail, then the Bengali for to roar, then say \emph{kā̃pāno} again
\end{itemize}

\begin{tcolorbox}[breakable,colback=teal!4,colframe=teal!35!black,title={Before you move on}]
What does \bn{কাঁপানো} mean? (\textquotedblleft{}To shake (something); to make tremble\textquotedblright{}.) Say it once more.
\end{tcolorbox}
\section[\texorpdfstring{\bn{নুয়ে} \bn{পড়া} (nuye pôṛā)}{নুয়ে পড়া (nuye pôṛā)}]{\bn{নুয়ে} \bn{পড়া} (nuye pôṛā) — to bend low; to droop}

\label{lesson:BN-C429-nuye-pora}



\begin{quote}
Before the new one: say the Bengali for to roar, then the Bengali for to shake (something).
\end{quote}

\subsection*{You'll want to know: \bn{নুয়ে} \bn{পড়া}}
\textbf{\bn{নুয়ে} \bn{পড়া}} — \emph{nuye pôṛā} — \textquotedblleft{}to bend low; to droop\textquotedblright{}.

\begin{quote}
\bn{ঝড়ে} \bn{গাছগুলো} \bn{নুয়ে} \bn{পড়ল।}
\end{quote}

\begin{quote}
\emph{jhôṛe gāchhgulo nuye pôṛlô.}
\end{quote}

\begin{quote}
The trees bent low in the storm.
\end{quote}

Past (he/she …): \bn{সেই} \bn{গাছ} \bn{নুয়ে} \bn{পড়ল}, \emph{sei gāchh nuye pôṛlô}.

Also of a person bowed with age or a load: \bn{বুড়ো} \bn{লোকটা} \bn{নুয়ে} \bn{হাঁটেন}.

\begin{cousinweb}
\textbf{\bn{নুয়ে} \bn{পড়া}} (\emph{nuye pôṛā}) joins \bn{নুয়ে}, the -\bn{এ} form of \bn{নোয়া}, \textquotedblleft{}to bow\textquotedblright{}, which CDIAL lists as Bengali \emph{nuyā}, \textquotedblleft{}bows, stoops\textquotedblright{}, under Sanskrit \emph{námati}, \textquotedblleft{}bends, bows\textquotedblright{} (CDIAL 6956), to \bn{পড়া}, to fall.
\end{cousinweb}

\begin{grammarlens}[title={What you've built}]
One more word for getting around.
\end{grammarlens}

\subsection*{Your turn}
\begin{itemize}
\raggedright
  \item \emph{Say it:} \emph{nuye pôṛā}
  \item \emph{Say it:} \emph{nuye pôṛā}, once more
  \item \emph{Say it:} \emph{jhôṛe gāchhgulo nuye pôṛlô.} — The trees bent low in the storm.
  \item \emph{Recall:} say the Bengali for to roar, then the Bengali for to shake (something), then say \emph{nuye pôṛā} again
\end{itemize}

\begin{tcolorbox}[breakable,colback=teal!4,colframe=teal!35!black,title={Before you move on}]
What does \bn{নুয়ে} \bn{পড়া} mean? (\textquotedblleft{}To bend low; to droop\textquotedblright{}.) Say it once more.
\end{tcolorbox}
